\documentclass[a4paper,11pt]{article}
\usepackage[a4paper,margin=0.68in,top=0.63in,bottom=0.60in]{geometry}
% Compile with XeLaTeX (Overleaf: Menu > Compiler > XeLaTeX).
\usepackage{fontspec}
\usepackage{xeCJK}
\setmainfont[Scale=0.97]{Liberation Sans}
\setsansfont[Scale=0.97]{Liberation Sans}
\setCJKmainfont{Noto Sans CJK TC}
\setCJKsansfont{Noto Sans CJK TC}
\renewcommand{\familydefault}{\sfdefault}
\usepackage{xcolor,tabularx,enumitem,titlesec,microtype,hyperref,fancyhdr}
\hypersetup{hidelinks}
\definecolor{ink}{HTML}{182333}
\definecolor{muted}{HTML}{536173}
\definecolor{rulecol}{HTML}{CBD3DF}
\definecolor{shade}{HTML}{F5F7FA}
\color{ink}
\setlength{\parindent}{0pt}
\setlength{\parskip}{0pt}
\linespread{1.055}
\setlength{\tabcolsep}{0pt}
\pagestyle{fancy}
\fancyhf{}
\fancyfoot[L]{\footnotesize\color{muted} Yu-Jui Liang}
\fancyfoot[R]{\footnotesize\color{muted} \thepage\ / 2}
\renewcommand{\headrulewidth}{0pt}\renewcommand{\footrulewidth}{0pt}
\setlength{\footskip}{17pt}
\newcommand{\dochead}[2]{%
  {\footnotesize\bfseries\color{muted} RESEARCH \& PROJECT SUMMARY}%
  \hfill{\footnotesize\color{muted} #1}\par\vspace{4pt}
  {\fontsize{20}{23}\selectfont\bfseries Yu-Jui Liang\hspace{0.4em}{\fontsize{17}{20}\selectfont 梁栯睿}}\par\vspace{3pt}
  {\small\color{muted}#2}\par\vspace{7pt}
  \color{rulecol}\rule{\linewidth}{1.1pt}\color{ink}\vspace{10pt}}
\newcommand{\smallhead}[1]{{\footnotesize\bfseries\color{muted} #1}\par\vspace{4pt}\color{rulecol}\rule{\linewidth}{0.8pt}\color{ink}\vspace{12pt}}
\newcommand{\metadata}[1]{{\small #1}\par\vspace{10pt}}
\newcommand{\project}[6]{%
{\small\color{muted}\bfseries #1}\par\vspace{3pt}
{\fontsize{12}{14}\selectfont\bfseries #2}\par\vspace{8pt}
{\small\textbf{Purpose.}\ #3}\par\vspace{6pt}
{\small\textbf{Approach.}\ #4}\par\vspace{6pt}
{\small\textbf{Evidence / Status.}\ #5}\par\vspace{6pt}
{\footnotesize\textbf{Relevant skills:}\ #6}\par\vspace{10pt}}
\newcommand{\flowline}[1]{\noindent\colorbox{shade}{\parbox{\dimexpr\linewidth-2\fboxsep}{\centering\footnotesize\color{muted} #1}}\par\vspace{18pt}}
\newcommand{\separate}{\par\noindent\color{rulecol}\rule{\linewidth}{0.7pt}\color{ink}\par\vspace{13pt}}
\begin{document}
\dochead{SOFTWARE / SYSTEMS}{Algorithms \quad / \quad System Simulation \quad / \quad Testing \& Validation}
\metadata{%
學歷：國立臺灣大學｜資料科學碩士學位學程（預計 2027 年 6 月畢業）\\[3pt]
\phantom{學歷：}國立政治大學｜資訊科學系學士（雙主修：電子物理學士學位學程）\\[3pt]
指導教授：江介宏教授（Jie-Hong Roland Jiang），國立臺灣大學\\[3pt]
共同指導：鄭湘筠博士（Hsiang-Yun Cheng），中央研究院\\[3pt]
專業領域：C++/Python, optimization, system simulation, GPU/memory systems, validation}
% ============================================================
% Project 01 - Heterogeneous GPU LLM Serving
% ============================================================

\project
{01 \quad / \quad GPU / MEMORY SYSTEMS (ONGOING)}
{Latency- and Energy-Aware Scheduling for Heterogeneous GPU LLM Serving}
{
Study efficient LLM inference serving across heterogeneous GPUs
with different compute capabilities, HBM / HBF memory configurations,
and prefill / decode execution characteristics.
The research focuses on resource imbalance, memory-access overhead,
and inter-device data movement as potential sources of inference
latency and energy inefficiency.
}
{
Designing a system-level simulation methodology to investigate
request-to-GPU scheduling, prefill / decode placement, memory
allocation, and data-movement policies under varying workloads.
The planned evaluation compares scheduling and memory-system
configurations using time to first token (TTFT), time per output
token (TPOT), end-to-end latency, throughput, power consumption,
and energy per generated token.
The objective is to identify configurations and policies that
improve latency and energy efficiency while respecting compute,
memory-capacity, and bandwidth constraints.
}
{
Ongoing research focused on developing a system-level
simulation framework for heterogeneous GPU LLM serving,
with emphasis on latency-aware scheduling, memory-system
optimization, and energy-efficient resource allocation.
}
{
C/C++, Python, Linux, Algorithm Design,
Optimization, System-Level Simulation,
Performance Modeling, Resource Scheduling,
Memory Management, Performance Analysis
}


\flowline{
GPU / Memory Configuration
$\longrightarrow$ Request Scheduling / Memory Policies
$\longrightarrow$ Latency / Energy Evaluation
}
\separate
\project{02 \quad / \quad LLM APPLICATIONS \& FORMAL REASONING}
{SAT-LEX: LLM and SMT-Based Document Verification System}
{Support more deterministic checking of domain-specific documents by combining extracted evidence with explicit rules and automated constraint solving.}
{Translated extracted document evidence and domain rules into structured constraints; applied an SMT solver to identify conflicting conditions and unsupported conclusions beyond unconstrained LLM responses.}
{The implemented constraint checks are intended to identify logical conflicts, inconsistent conditions, and unsupported conclusions.}
{LLM applications, SMT / constraint solving, automated reasoning, verification workflows}
\flowline{Rules / Document Evidence $\longrightarrow$ Structured Constraints $\longrightarrow$ Solver Checks}
\newpage
\smallhead{SELECTED PROJECTS \quad | \quad SOFTWARE / SYSTEMS}
\project{03 \quad / \quad BACKEND SYSTEMS \& SOFTWARE VALIDATION}
{EV Charging Management System -- NeoPower Technologies}
{Develop and validate an electric-bus charging management platform integrating vehicles, charging equipment, and external vendor services, with reliable operation across customer sites.}
{Contributed to system planning and backend development, including APIs, database integration, and log-based debugging. Used a container-based architecture to isolate projects and services on shared servers, reducing interference during deployment and troubleshooting. Independently developed a custom simulator to emulate external vendors acting as clients or servers. Reconstructed normal and failure scenarios from historical database logs and customer feedback, including sudden charging-demand surges and simultaneous charger disconnections caused by site-network failures.}
{Professional engineering experience (Jun. 2023 -- Mar. 2025). Used scenario-based simulation and regression testing before deployment, followed by validation using production logs, operational data, and customer feedback. Supported A/B-style comparisons of existing and revised algorithms across sites with similar conditions, such as charging-power allocation policies, to assess behavior under real operating conditions.}
{Backend development, API / database integration, containers, vendor simulation, regression testing, production validation, log analysis, debugging}
\flowline{Vendor Simulation $\longrightarrow$ Scenario / Regression Tests $\longrightarrow$ Deployment $\longrightarrow$ Operational Validation}
\separate
\project{04 \quad / \quad LOGIC SYNTHESIS \& TECHNOLOGY MAPPING}
{Multi-Output Standard-Cell Mapping with MILP and Differentiable Optimization}
{Improve area--delay product (ADP) by exploiting multi-output standard cells in the academic ASAP7 7-nm technology library, compared with conventional single-output mapping.}
{Extended the ABC logic-synthesis tool in C++ with a custom multi-output (MO) cut-discovery path and Boolean matching. Applied area-aware mixed-integer linear programming (MILP) to select matched cut/cell pairs, then passed the selected candidates to a PyTorch-based differentiable mapping optimizer for Liberty-based area/delay evaluation.}
{Validated correctness and evaluated 30 arithmetic benchmarks. The updated MILP-enhanced mapping flow achieved a 6\% average ADP improvement over the single-output baseline.}
{C++, ABC, logic synthesis, technology mapping, Boolean matching, PyTorch, Liberty, ASAP7, QoR analysis}
\flowline{ABC / MO Cut Discovery $\longrightarrow$ Boolean Matching $\longrightarrow$ Area-Aware MILP Selection $\longrightarrow$ Differentiable Mapping}
\vfill
\end{document}
